
Or @@M0@@ est divisible par @@M1@@ :
pour le voir, on peut se placer dans le @@M2@@-anneau universel qui sert à définir les @@M3@@ (V 5); alors @@M4@@ est un polynôme en @@M5@@, nul pour @@M6@@, donc divisible par @@M7@@. Par ailleurs, l'équation de dépendance de @@M8@@ sur @@M9@@ est @@M10@@ (VI (1.13.1)). On obtient donc :
@@M11@@
En appliquant @@M12@@, on obtient :
@@M13@@
car @@M14@@.

D'après la formule de (V 5.6), appliquée à @@M15@@, @@M16@@ est le coefficient de @@M17@@ dans la série :
@@M18@@

De même, la même formule appliquée à @@M19@@ et @@M20@@ montre que @@M21@@ est le coefficient de @@M22@@ dans la série :
@@M23@@

Les congruences (4.3.4) peuvent donc s'écrire :
@@M24@@
